\chapter{Soundness and Completeness of Propositional Logic}\label{ch:soundcomplete}

We now have two notions of consequence: $\Sigma\models\varphi$ (semantic, via truth
valuations) and $\Sigma\vdash\varphi$ (syntactic, via natural deduction). This chapter
proves they coincide.

\begin{important}[The two theorems]
\begin{description}
  \item[Soundness:] if $\Sigma\vdash\varphi$ then $\Sigma\models\varphi$.
        (\emph{Everything provable is true}: the rules never lead us astray.)
  \item[Completeness:] if $\Sigma\models\varphi$ then $\Sigma\vdash\varphi$.
        (\emph{Everything true is provable}: the rules are strong enough.)
\end{description}
\end{important}

\section{Why we care}

\begin{itemize}
  \item Soundness gives a way to show something is \emph{not} provable: if we find a
        valuation making $\Sigma$ true and $\varphi$ false, then $\Sigma\not\models\varphi$,
        so $\Sigma\nvdash\varphi$. Proving non-provability directly (``no proof
        exists'') would otherwise be very hard.
  \item Completeness says that whenever a truth table confirms an entailment, a formal
        proof exists, so we may search for one with confidence.
\end{itemize}

\section{Soundness}

\begin{theorem}[Soundness]
If $\varphi_1,\dots,\varphi_n\vdash\psi$ is valid (has a natural deduction proof) then
$\varphi_1,\dots,\varphi_n\models\psi$.
\end{theorem}

\begin{proof}
We prove a stronger statement by (strong) induction on the line number $k$ of a proof:

\emph{$M(k)$: if the formula on line $k$ is $\chi$, and $\alpha_1,\dots,\alpha_m$ are the
assumptions of the boxes open at line $k$, then
$\varphi_1,\dots,\varphi_n,\alpha_1,\dots,\alpha_m\models\chi$.}

Applying $M(k)$ to the last line (where no boxes are open) gives the theorem. Assume
$M(j)$ for all $j<k$ and consider the justification of line $k$.
\begin{itemize}
  \item \emph{Premise}: $\chi=\varphi_i$, and trivially $\varphi_1,\dots,\varphi_n\models\varphi_i$.
  \item \emph{Assumption}: $\chi=\alpha_m$ is itself one of the open assumptions, so the
        entailment is trivial.
  \item \emph{$\land i$} from lines $j_1,j_2<k$ with $\chi=\chi_1\land\chi_2$: lines $j_1,j_2$
        are visible from $k$, so their open assumptions are among those at $k$. By IH any
        valuation satisfying the premises and open assumptions makes $\chi_1$ and
        $\chi_2$ true, hence $\chi$ true. (Monotonicity handles the extra assumptions.)
  \item \emph{$\to e$} from $\chi_1$ and $\chi_1\to\chi$: any valuation making both true
        makes $\chi$ true.
  \item \emph{$\to i$} with $\chi=\chi_1\to\chi_2$ from the box $i$--$j$: by IH applied to
        line $j$, $\varphi$'s, $\alpha$'s, $\chi_1\models\chi_2$. By the semantic
        deduction theorem, $\varphi$'s, $\alpha$'s $\models\chi_1\to\chi_2$.
  \item \emph{$\neg i$} from a box $\chi_1\ \dots\ \bot$: by IH the premises, open
        assumptions and $\chi_1$ entail $\bot$, i.e.\ no valuation satisfies them all; so
        every valuation satisfying the premises and open assumptions makes $\chi_1$ false,
        and $\neg\chi_1$ true.
  \item \emph{$\lor e$}: let $t$ satisfy the premises and open assumptions. By IH
        $(\chi_1\lor\chi_2)^t=\TT$. If $\chi_1^t=\TT$, the IH on the last line of the first
        box gives $\chi^t=\TT$; otherwise $\chi_2^t=\TT$ and the second box gives it.
  \item The remaining rules ($\land e$, $\lor i$, $\neg e$, $\bot e$, $\neg\neg e$) are
        similar: in each case the conclusion is true under every valuation that makes
        the cited lines true. ($\bot e$ holds vacuously: no valuation makes $\bot$ true.)
\end{itemize}
\end{proof}

\begin{example}[Using soundness]
Is $p\to q\vdash q\to p$ valid? Take $t(p)=\FF$, $t(q)=\TT$: then $(p\to q)^t=\TT$ but
$(q\to p)^t=\FF$. So $p\to q\not\models q\to p$, and by soundness there is no proof.
\end{example}

\section{Completeness}

We follow the three-step proof from Huth and Ryan for finite premise sets, then outline
the general (maximal consistent set) proof.

\subsection*{Step 1: reduce to theorems}
If $\varphi_1,\dots,\varphi_n\models\psi$, then
\[
  \models \varphi_1\to(\varphi_2\to(\dots(\varphi_n\to\psi)\dots)).
\]
(If a valuation makes this false, it makes all $\varphi_i$ true and $\psi$ false.)

\subsection*{Step 2: every tautology is a theorem}
\begin{lemma}[Kalm\'ar]
Let $\varphi$ have atoms $p_1,\dots,p_n$, and let $t$ be a valuation. Define
$\hat p_i = p_i$ if $t(p_i)=\TT$ and $\hat p_i=\neg p_i$ otherwise. Then
\[
  \hat p_1,\dots,\hat p_n\vdash\varphi \text{ if } \varphi^t=\TT,\qquad
  \hat p_1,\dots,\hat p_n\vdash\neg\varphi \text{ if } \varphi^t=\FF.
\]
\end{lemma}
\begin{proof}[Proof sketch]
Structural induction on $\varphi$. Base: $\varphi=p_i$, and $\hat p_i$ is a premise. Step
$\neg\varphi_1$: if $\varphi_1^t=\TT$ then $(\neg\varphi_1)^t=\FF$ and we need
$\neg\neg\varphi_1$, obtained from $\varphi_1$ by $\neg\neg i$; otherwise IH gives
$\neg\varphi_1$ directly. Step $\varphi_1\to\varphi_2$: three cases. If $\varphi_2^t=\TT$,
derive $\varphi_2$ by IH and then $\varphi_1\to\varphi_2$ by $\to i$ (with a box that just
copies $\varphi_2$). If $\varphi_1^t=\FF$, derive $\neg\varphi_1$, then inside a box assume
$\varphi_1$, get $\bot$, then $\varphi_2$ by $\bot e$. If $\varphi_1^t=\TT,\varphi_2^t=\FF$,
derive $\varphi_1$ and $\neg\varphi_2$ and then $\neg(\varphi_1\to\varphi_2)$ (assume the
implication, apply $\to e$, get $\bot$). The other connectives are similar.
\end{proof}

Now suppose $\models\eta$ with atoms $p_1,\dots,p_n$. By the lemma, for each of the
$2^n$ valuations, $\hat p_1,\dots,\hat p_n\vdash\eta$. Use LEM on $p_1$:
$p_1\lor\neg p_1$, and apply $\lor e$. In each case, use LEM on $p_2$, and so on. At the
leaves of this nested case analysis all of $\hat p_1,\dots,\hat p_n$ are available, and we
plug in the proof from the lemma. Altogether $\vdash\eta$.

\subsection*{Step 3: put the premises back}
From $\vdash\varphi_1\to(\dots\to\psi)$, add premises $\varphi_1,\dots,\varphi_n$ and use
$\to e$ $n$ times to obtain $\varphi_1,\dots,\varphi_n\vdash\psi$. \qed

\subsection{The general completeness theorem}

For possibly infinite $\Sigma$ we use consistency.

\begin{definition}
$\Sigma$ is \emph{consistent} if there is no $\varphi$ with $\Sigma\vdash\varphi$ and
$\Sigma\vdash\neg\varphi$ (equivalently, $\Sigma\nvdash\bot$). $\Sigma$ is
\emph{maximal consistent} if it is consistent and for every formula $\varphi$, either
$\varphi\in\Sigma$ or $\neg\varphi\in\Sigma$.
\end{definition}

\begin{lemma}
$\Sigma\nvdash\varphi$ iff $\Sigma\cup\{\neg\varphi\}$ is consistent.
\end{lemma}

\begin{theorem}[Lindenbaum's lemma]
Every consistent set extends to a maximal consistent set.
\end{theorem}
\begin{proof}[Sketch]
Enumerate all formulas $\varphi_1,\varphi_2,\dots$ (possible because $\Prop$ is countable).
Let $\Sigma_0=\Sigma$ and $\Sigma_{k+1}=\Sigma_k\cup\{\varphi_{k+1}\}$ if that is
consistent, else $\Sigma_k\cup\{\neg\varphi_{k+1}\}$. One shows each $\Sigma_k$ is
consistent, and so is $\Sigma^*=\bigcup_k\Sigma_k$ because proofs are finite.
\end{proof}

\begin{theorem}
Every consistent set is satisfiable.
\end{theorem}
\begin{proof}[Sketch]
Extend to a maximal consistent $\Sigma^*$ and define $t(p)=\TT$ iff $p\in\Sigma^*$. By
structural induction, $\varphi^t=\TT$ iff $\varphi\in\Sigma^*$ for every formula $\varphi$
(e.g.\ $(\neg\varphi)\in\Sigma^*$ iff $\varphi\notin\Sigma^*$ by maximality and
consistency). So $t$ satisfies $\Sigma\subseteq\Sigma^*$.
\end{proof}

\begin{corollary}[Completeness]
If $\Sigma\models\varphi$ then $\Sigma\vdash\varphi$.
\end{corollary}
\begin{proof}
If $\Sigma\nvdash\varphi$, then $\Sigma\cup\{\neg\varphi\}$ is consistent, hence satisfiable;
a valuation satisfying it shows $\Sigma\not\models\varphi$.
\end{proof}

\section{Compactness}

\begin{theorem}[Compactness]
A set $\Sigma$ of formulas is satisfiable iff every finite subset of $\Sigma$ is
satisfiable. Equivalently, $\Sigma\models\varphi$ iff $\Sigma_0\models\varphi$ for some
finite $\Sigma_0\subseteq\Sigma$.
\end{theorem}
\begin{proof}
If every finite subset is satisfiable, then (by soundness) every finite subset is
consistent. Since a derivation of $\bot$ uses only finitely many premises, $\Sigma$ is
consistent, and by the theorem above, satisfiable.
\end{proof}

Compactness is a purely semantic fact proved here through the proof system: a nice
illustration of why soundness and completeness are useful together.

\section{Decidability}

For finite $\Sigma$, $\Sigma\models\varphi$ is \emph{decidable}: check all $2^n$ rows of the
truth table. By completeness, provability $\Sigma\vdash\varphi$ is therefore decidable too.
A polynomial-time algorithm for tautology checking exists iff P $=$ NP, which is
open (SAT is NP-complete; TAUT is co-NP-complete).
We will see in Chapter~\ref{ch:undecidability} that the analogous problem for predicate
logic is undecidable.

\section*{Summary}
\begin{itemize}
  \item Soundness ($\vdash\Rightarrow\models$): induction on proof lines, checking each
        rule preserves truth relative to the premises and open assumptions.
  \item Use soundness plus a counterexample valuation to show a sequent has no proof.
  \item Completeness ($\models\Rightarrow\vdash$): finite case via Kalm\'ar's lemma and
        LEM case splits; general case via consistent $\Rightarrow$ maximal consistent
        $\Rightarrow$ satisfiable.
  \item Compactness: satisfiable iff every finite subset is satisfiable.
  \item Propositional validity is decidable (by truth tables), though possibly
        exponentially expensive.
\end{itemize}
